% GENERATED FILE. Edit canonical lessons, then run npm run generate:books.
% canonical-source-hash: fnv1a64:3c7bd1d0df4f63d1
% canonical-lessons: BN-C364-kothin, BN-C364-songkot, BN-C364-daridryo, BN-C364-biporjosto-hooya, BN-C364-chorom

\chapter{Hard Times}
\label{ch:bn-hard-times}
\hlchaptermodality{364}

\begin{chapteropening}
\textbf{By the end of this chapter:} \emph{I can talk about a crisis and how people got through it.}

Complete the last lesson of chapter 364: I can talk about a crisis and how people got through it.
\end{chapteropening}

\section[\texorpdfstring{\bn{কঠিন} (kôṭhin)}{কঠিন (kôṭhin)}]{\bn{কঠিন} (kôṭhin) — hard; difficult}

\label{lesson:BN-C364-kothin}



\begin{quote}
Before the new one: say the Bengali for an answer, then the Bengali for to come through.
\end{quote}

\subsection*{You'll want to know: \bn{কঠিন}}
\textbf{\bn{কঠিন}} — \emph{kôṭhin} — \textquotedblleft{}hard; difficult\textquotedblright{}.

\begin{quote}
\bn{পরীক্ষাটা} \bn{কঠিন} \bn{ছিল।}
\end{quote}

\begin{quote}
\emph{pôrīkkhāṭā kôṭhin chhilô.}
\end{quote}

\begin{quote}
The exam was hard.
\end{quote}

\begin{cousinweb}
\textbf{\bn{কঠিন}} (\emph{kôṭhin}) is Sanskrit \emph{kaṭhina}, hard or firm (Monier-Williams). It is hard to the touch and hard to do; \bn{সহজ}, easy, is its opposite.
\end{cousinweb}

\begin{grammarlens}[title={What you've built}]
One more word for giving reasons and opinions.
\end{grammarlens}

\subsection*{Your turn}
\begin{itemize}
\raggedright
  \item \emph{Say it:} \emph{kôṭhin}
  \item \emph{Say it:} \emph{kôṭhin}, once more
  \item \emph{Say it:} \emph{pôrīkkhāṭā kôṭhin chhilô.} — The exam was hard.
  \item \emph{Recall:} say the Bengali for an answer, then the Bengali for to come through, then say \emph{kôṭhin} again
\end{itemize}

\begin{tcolorbox}[breakable,colback=teal!4,colframe=teal!35!black,title={Before you move on}]
What does \bn{কঠিন} mean? (\textquotedblleft{}Hard; difficult\textquotedblright{}.) Say it once more.
\end{tcolorbox}
\section[\texorpdfstring{\bn{সংকট} (sôngkôṭ)}{সংকট (sôngkôṭ)}]{\bn{সংকট} (sôngkôṭ) — a crisis}

\label{lesson:BN-C364-songkot}



\begin{quote}
Before the new one: say the Bengali for to come through, then the Bengali for difficult.
\end{quote}

\subsection*{You'll want to know: \bn{সংকট}}
\textbf{\bn{সংকট}} — \emph{sôngkôṭ} — \textquotedblleft{}a crisis\textquotedblright{}.

\begin{quote}
\bn{দেশে} \bn{এখন} \bn{জলের} \bn{সংকট।}
\end{quote}

\begin{quote}
\emph{deshe ekhôn jôler sôngkôṭ.}
\end{quote}

\begin{quote}
There is a water crisis in the country now.
\end{quote}

\begin{cousinweb}
\textbf{\bn{সংকট}} (\emph{sôngkôṭ}) is Sanskrit \emph{saṃkaṭa}, narrow and so dangerous or critical (Monier-Williams). In Bengali it is the noun, a crisis or a severe shortage.
\end{cousinweb}

\begin{grammarlens}[title={What you've built}]
One more word for giving reasons and opinions.
\end{grammarlens}

\subsection*{Your turn}
\begin{itemize}
\raggedright
  \item \emph{Say it:} \emph{sôngkôṭ}
  \item \emph{Say it:} \emph{sôngkôṭ}, once more
  \item \emph{Say it:} \emph{deshe ekhôn jôler sôngkôṭ.} — There is a water crisis in the country now.
  \item \emph{Recall:} say the Bengali for to come through, then the Bengali for difficult, then say \emph{sôngkôṭ} again
\end{itemize}

\begin{tcolorbox}[breakable,colback=teal!4,colframe=teal!35!black,title={Before you move on}]
What does \bn{সংকট} mean? (\textquotedblleft{}A crisis\textquotedblright{}.) Say it once more.
\end{tcolorbox}
\section[\texorpdfstring{\bn{দারিদ্র্য} (dāridryô)}{দারিদ্র্য (dāridryô)}]{\bn{দারিদ্র্য} (dāridryô) — poverty}

\label{lesson:BN-C364-daridryo}



\begin{quote}
Before the new one: say the Bengali for difficult, then the Bengali for a crisis.
\end{quote}

\subsection*{You'll want to know: \bn{দারিদ্র্য}}
\textbf{\bn{দারিদ্র্য}} — \emph{dāridryô} — \textquotedblleft{}poverty\textquotedblright{}.

\begin{quote}
\bn{বন্যার} \bn{পরে} \bn{গ্রামে} \bn{দারিদ্র্য} \bn{দেখা} \bn{দিল।}
\end{quote}

\begin{quote}
\emph{bônyār pôre grāme dāridryô dækhā dilô.}
\end{quote}

\begin{quote}
Poverty took hold in the village after the flood.
\end{quote}

\begin{cousinweb}
\textbf{\bn{দারিদ্র্য}} (\emph{dāridryô}) is Sanskrit \emph{dāridrya}, poverty (Monier-Williams), the abstract of \emph{daridra} 'poor'. The everyday adjective is the taught \bn{গরিব}.
\end{cousinweb}

\begin{grammarlens}[title={What you've built}]
One more word for giving reasons and opinions.
\end{grammarlens}

\subsection*{Your turn}
\begin{itemize}
\raggedright
  \item \emph{Say it:} \emph{dāridryô}
  \item \emph{Say it:} \emph{dāridryô}, once more
  \item \emph{Say it:} \emph{bônyār pôre grāme dāridryô dækhā dilô.} — Poverty took hold in the village after the flood.
  \item \emph{Recall:} say the Bengali for difficult, then the Bengali for a crisis, then say \emph{dāridryô} again
\end{itemize}

\begin{tcolorbox}[breakable,colback=teal!4,colframe=teal!35!black,title={Before you move on}]
What does \bn{দারিদ্র্য} mean? (\textquotedblleft{}Poverty\textquotedblright{}.) Say it once more.
\end{tcolorbox}
\section[\texorpdfstring{\bn{বিপর্যস্ত} \bn{হওয়া} (bipôrjôstô hôoyā)}{বিপর্যস্ত হওয়া (bipôrjôstô hôoyā)}]{\bn{বিপর্যস্ত} \bn{হওয়া} (bipôrjôstô hôoyā) — to be devastated; to be thrown into turmoil}

\label{lesson:BN-C364-biporjosto-hooya}



\begin{quote}
Before the new one: say the Bengali for a crisis, then the Bengali for poverty.
\end{quote}

\subsection*{You'll want to know: \bn{বিপর্যস্ত} \bn{হওয়া}}
\textbf{\bn{বিপর্যস্ত} \bn{হওয়া}} — \emph{bipôrjôstô hôoyā} — \textquotedblleft{}to be devastated; to be thrown into turmoil\textquotedblright{}.

\begin{quote}
\bn{বন্যায়} \bn{পুরো} \bn{গ্রাম} \bn{বিপর্যস্ত} \bn{হল।}
\end{quote}

\begin{quote}
\emph{bônyāy puro grām bipôrjôstô hôlô.}
\end{quote}

\begin{quote}
The whole village was devastated by the flood.
\end{quote}

Past (he/she …): \bn{সেই} \bn{শহরটা} \bn{ঝড়ে} \bn{বিপর্যস্ত} \bn{হল}, \emph{sei shôhôrṭā jhôṛe bipôrjôstô hôlô}.

\begin{cousinweb}
\textbf{\bn{বিপর্যস্ত} \bn{হওয়া}} (\emph{bipôrjôstô hôoyā}) is Sanskrit \emph{viparyasta}, turned over or upset (Monier-Williams), with \bn{হওয়া}: a town or a life turned upside down by a disaster. The sense 'devastated' is a Bengali use of that picture; the cause takes -\bn{এ}: \bn{বন্যায়}.
\end{cousinweb}

\begin{grammarlens}[title={What you've built}]
One more word for giving reasons and opinions.
\end{grammarlens}

\subsection*{Your turn}
\begin{itemize}
\raggedright
  \item \emph{Say it:} \emph{bipôrjôstô hôoyā}
  \item \emph{Say it:} \emph{bipôrjôstô hôoyā}, once more
  \item \emph{Say it:} \emph{bônyāy puro grām bipôrjôstô hôlô.} — The whole village was devastated by the flood.
  \item \emph{Recall:} say the Bengali for a crisis, then the Bengali for poverty, then say \emph{bipôrjôstô hôoyā} again
\end{itemize}

\begin{tcolorbox}[breakable,colback=teal!4,colframe=teal!35!black,title={Before you move on}]
What does \bn{বিপর্যস্ত} \bn{হওয়া} mean? (\textquotedblleft{}To be devastated; to be thrown into turmoil\textquotedblright{}.) Say it once more.
\end{tcolorbox}
\section[\texorpdfstring{\bn{চরম} (chôrôm)}{চরম (chôrôm)}]{\bn{চরম} (chôrôm) — extreme; utmost}

\label{lesson:BN-C364-chorom}



\begin{quote}
Before the new one: say the Bengali for poverty, then the Bengali for to be devastated.
\end{quote}

\subsection*{You'll want to know: \bn{চরম}}
\textbf{\bn{চরম}} — \emph{chôrôm} — \textquotedblleft{}extreme; utmost\textquotedblright{}.

\begin{quote}
\bn{বন্যার} \bn{পরে} \bn{গ্রামের} \bn{মানুষ} \bn{চরম} \bn{বিপদে} \bn{পড়ল।}
\end{quote}

\begin{quote}
\emph{bônyār pôre grāmer mānush chôrôm bipôde pôṛlô.}
\end{quote}

\begin{quote}
After the flood the villagers were in dire trouble.
\end{quote}

\begin{cousinweb}
\textbf{\bn{চরম}} (\emph{chôrôm}) is Sanskrit \emph{carama}, last or ultimate (Monier-Williams): at the far end of the scale. Bengali uses it for the worst a hardship can get, \bn{চরম} \bn{বিপদ}, or the height of anything, \bn{চরম} \bn{আনন্দ}; that use is not one the dictionary gives.
\end{cousinweb}

\begin{grammarlens}[title={What you've built}]
One more word for giving reasons and opinions.
\end{grammarlens}

\subsection*{Your turn}
\begin{itemize}
\raggedright
  \item \emph{Say it:} \emph{chôrôm}
  \item \emph{Say it:} \emph{chôrôm}, once more
  \item \emph{Say it:} \emph{bônyār pôre grāmer mānush chôrôm bipôde pôṛlô.} — After the flood the villagers were in dire trouble.
  \item \emph{Recall:} say the Bengali for poverty, then the Bengali for to be devastated, then say \emph{chôrôm} again
\end{itemize}

\begin{tcolorbox}[breakable,colback=teal!4,colframe=teal!35!black,title={Before you move on}]
What does \bn{চরম} mean? (\textquotedblleft{}Extreme; utmost\textquotedblright{}.) Say it once more.
\end{tcolorbox}
